% paper/sections/00_abstract.tex
\begin{titlepage}
\begin{center}
\vspace*{12mm}
{\color{navy}\rule{\linewidth}{2pt}}\\[8pt]
{\LARGE\bfseries 高次コンパクト差分に基づく気液二相流数値法}\\[6pt]
{\large\bfseries —— 保存形界面追跡・Balanced-Force・圧力ジャンプ分相 PPE の統合設計 ——}\\[4pt]
{\color{navy}\rule{\linewidth}{2pt}}\\[10pt]
{\large
\textbf{High-Order Compact Schemes for Gas--Liquid Two-Phase Flow Simulation}\\[4pt]
\textit{Conservative Level Set $\cdot$ Balanced Force $\cdot$ Pressure-Jump Phase-Separated PPE}
}
\vspace{14mm}
\begin{tcolorbox}[mybox, width=0.88\textwidth]
\small\centering
\begin{tabular}{@{}p{0.9\linewidth}@{}}
\textbf{本稿が答える問い}\\[3pt]
気液界面を何として保存し，どの面位置で圧力・毛管力・速度補正を閉じると，
寄生流れと体積ドリフトを同時に抑えられるか．\\[6pt]
\textbf{中心の読み筋}\\[3pt]
保存形界面量を高次面演算子，分相 PPE，HFE，Balanced-Force 条件へ
同じ面位置で渡し，単体検証・統合検証・物理ベンチマークを通じて
成立範囲と未達課題を分けて示す．
\end{tabular}
\end{tcolorbox}
\end{center}
\vfill
{\color{navy}\rule{\linewidth}{0.8pt}}\\[2pt]
{\small\color{darkgray}対象読者：大学院生・研究者（計算流体力学の基礎知識を前提）}
\end{titlepage}

% ─────────────────────────────────────────────────────────────────────────────
\begin{abstract}
気液二相流数値シミュレーションにおける主要な数値的障害は，
静止すべき界面に生じる人工的渦流れ——\textbf{寄生流れ（spurious currents）}——である．
本稿は，その離散化起因成分を構造的に抑えるため，
結合コンパクト差分法（CCD，内点 $\Ord{h^6}$），FCCD 面ジェット，UCCD6，
Conservative Level Set（CLS）法，Hermite 場延長法（HFE；\cite{Aslam2004} の Extension PDE に基づく），
および Balanced-Force 七原則を一つの圧力ジャンプ型分相 PPE 構成へ統合する．
標準経路では，界面を FCCD 保存形 CLS または幾何セル分率 $q_C$ の保存輸送で更新し，
表面張力は CSF 体積力ではなく，Young--Laplace 束仮想仕事と
$j_{gl}=p_g-p_l=-\sigma\kappa_{lg}$ の圧力ジャンプとして
分相 FCCD-PPE + HFE + 欠陥補正（DC 収束判定）へ組み込む．
この構成により，コロケート格子でも Rhie--Chow 補間に依存せず，
圧力勾配・表面張力・HFE 上流点・GFM ジャンプが同じ面位置とステンシルを共有する
Balanced-Force 部分系として閉じる．
時間発展は TVD-RK3（CLS），IMEX-BDF2/EXT2（運動量対流），
implicit-BDF2 + DC（粘性），IPC 圧力増分を組み合わせた 7 ステップ手順として整理する．
移動格子のアクティブ幾何毛管分解では，射影後の面フラックス量を
格子再構築後も同じ仕事量として扱い，圧力履歴は滑らかな座標だけを外挿し，
現在の幾何毛管反力は履歴へ混入しない．

検証は，単体検証 U1--U12，統合検証 V1--V12，および
物理ベンチマーク第~\ref{sec:validation}章の三層で構成した．
単体検証では CCD $\Ord{h^{6.0}}$，FCCD $\Ord{h^{4.00}}$，
UCCD6 $\Ord{h^{7.00}}$，HFE 1D $\Ord{h^{5.91}}$，implicit-BDF2 $\Ord{\Delta t^{2.01}}$，
保存形共通フラックス，制約付き面速度状態・周期壁商空間条件，
アクティブ幾何毛管分解の採用境界などを確認し，
統合検証では静止液滴 Laplace 圧力誤差 $0.97\%$，
寄生流れの CCD/FD 比 $1/2.5$--$1/5.85$（代表ケース），
密度比 $\rho_l/\rho_g\le833$ での圧力ジャンプ型分相 PPE + DC + HFE のロバスト性
（$u_\infty^{\mathrm{end}}\le1.3\times10^{-9}$，$|\Delta V_\psi|/V_{\psi,0}\le1.2\times10^{-16}$），
Zalesak / 単一渦の非連成 CLS 強変形移流の質量・重心診断，
および採用条件を満たさない毛管分解を標準構成へ混入させない V11 ゲートを示した．
一方，圧力ジャンプを含む二相結合構成の時間精度は V7 で最終収束率 $1.59$ に留まり，
毛管圧力ジャンプと射影の界面帯誤差に支配される条件付き合格として扱う．
さらに物理ベンチマークでは，PhaseRegion 縮約診断で読む毛管波の毛管復元力と
Rayleigh--Lamb 型振動液滴の短時間復元，
標準実行で残る 1 周期の鋭い界面体積保存課題，
単一気泡上昇の浮力--投影閉包，
および Rayleigh--Taylor 線形成長から非線形形状形成までを，
採用済み経路と未接続経路を分けて整理した．
\end{abstract}
